% Contenidos del capítulo.
% Las secciones presentadas son orientativas y no representan
% necesariamente la organización que debe tener este capítulo.

\section{Introducción}

La inversión en mercados financieros ha ganado popularidad en los últimos años, facilitada por el acceso a plataformas digitales. Sin embargo, la falta de educación financiera lleva a muchos inversores novatos a cometer errores costosos. Según estudios clásicos de Barber y Odean, basados en más de 60.000 hogares estadounidenses, los inversores particulares que operan con mayor frecuencia obtienen rentabilidades anuales entre 5 y 10 puntos porcentuales inferiores a las de un índice de mercado amplio, principalmente debido a sesgos conductuales como el exceso de confianza y a una operativa excesiva \cite{barber2000trading}.

Este TFG presenta el diseño e implementación de una aplicación web orientada al análisis fundamental de empresas y a la simulación educativa de estrategias de inversión a largo plazo. La herramienta permite practicar estrategias como el \textit{Dollar Cost Averaging} (DCA) en un entorno libre de riesgo, utilizando datos históricos reales del mercado obtenidos a través de Yahoo Finance \cite{yfinance}. En su estado final, la aplicación incorpora además autenticación por correo, modo invitado, persistencia selectiva, historial para usuario registrado y varios módulos complementarios de apoyo pedagógico.

A diferencia de los simuladores de \textit{trading} tradicionales, que suelen fomentar la especulación a corto plazo, esta propuesta se centra en la inversión sosegada y el análisis de métricas fundamentales (PER, deuda, crecimiento), integrando un modo carrera gamificado para evaluar el desempeño del usuario ante distintos escenarios macroeconómicos y un ecosistema de uso más amplio compuesto por secciones de aprendizaje, validación progresiva y análisis complementario.

\section{Estructura de la memoria}

El resto de este documento se estructura de la siguiente manera:

\begin{itemize}
    \item En el \textbf{Capítulo 2} se detallan la motivación del proyecto y los objetivos principales y específicos que se pretenden alcanzar.
    \item El \textbf{Capítulo 3} revisa el estado del arte, analizando las herramientas existentes y justificando las tecnologías seleccionadas.
    \item El \textbf{Capítulo 4} aborda la planificación del proyecto, incluyendo la especificación de requisitos, la evolución metodológica, la viabilidad y los riesgos identificados.
    \item El \textbf{Capítulo 5} describe el desarrollo del sistema, desde la arquitectura funcional hasta la implementación de los módulos principales y el rediseño frontend.
    \item El \textbf{Capítulo 6} presenta las pruebas realizadas y los resultados obtenidos.
    \item El \textbf{Capítulo 7} expone las conclusiones y las líneas de trabajo futuro.
    \item Los \textbf{apéndices} recogen el manual técnico y de despliegue, así como un breve manual de usuario.
\end{itemize}